% ------------------------------------ Forum Acusticum 2026 Migrated from ICASSP 2024 submission
% ------------------------------------
\documentclass[non-peer-review]{fa2026}

% Additional packages (amsmath, amssymb, graphicx already loaded by fa2026.cls)
\usepackage{xparse,subcaption,acronym} \usepackage{amsthm, mathdots, mathtools} \newlength{\hsvmapheight}
\acrodef{CAP}{common-acoustical-pole}
\acrodef{CAPZ}{common-acoustical-pole and zero}
\acrodef{ERA}{Eigensystem Realization Algorithm}
\acrodef{FIR}{finite impulse response}
\acrodef{HATO}{head-above-torso orientation}
\acrodef{HRIR}{head-related impulse response}
\acrodef{HRTF}{head-related transfer function}
\acrodef{ITD}{interaural time delay}
\acrodef{LTI}{linear time-invariant}
\acrodef{MIMO}{multiple-input multiple-output}
\acrodef{MISO}{multiple-input single-output}
\acrodef{SIMO}{single-input multiple-output}
\acrodef{SVD}{singular value decomposition}
\acrodef{TERA}{Tangential Eigensystem Realization Algorithm}
\acrodef{UPOLS}{uniformly partitioned overlap-save}

\acrodefplural{HATO}[HATOs]{head-above-torso orientations}
\acrodefplural{HRIR}[HRIRs]{head-related impulse responses}
\acrodefplural{HRTF}[HRTFs]{head-related transfer functions}
\acrodefplural{ITD}[ITDs]{interaural time delays}
 \input{../generated/data/experiment-values} \mathtoolsset{mathic=true}
\newtheorem{theorem}{Theorem}[section] \crefformat{equation}{(#2#1#3)} \crefmultiformat{equation}{(#2#1#3)}{,
  (#2#1#3)}{, (#2#1#3)}{, and (#2#1#3)} \crefrangeformat{equation}{(#3#1#4--#5#2#6)}

% Math notation commands
\NewDocumentCommand{\N}{oo}{\ensuremath{\mathbb{N}\IfValueTF{#1}{\IfValueTF{#2}{^{#1\times#2}}{^#1}}{}}}
\NewDocumentCommand{\R}{oo}{\ensuremath{\mathbb{R}\IfValueTF{#1}{\IfValueTF{#2}{^{#1\times#2}}{^#1}}{}}}
\NewDocumentCommand{\C}{oo}{\ensuremath{\mathbb{C}\IfValueTF{#1}{\IfValueTF{#2}{^{#1\times#2}}{^#1}}{}}}
\NewDocumentCommand{\Hardy}{o}{\ensuremath{\mathcal{H}_{#1}}}
\NewDocumentCommand{\rank}{m}{\ensuremath{\operatorname{rank}\left(#1\right)}}
\NewDocumentCommand{\diag}{m}{\ensuremath{\operatorname{diag}\left(#1\right)}}
\NewDocumentCommand{\blkdiag}{m}{\ensuremath{\operatorname{blk\,diag}\left(#1\right)}}
\NewDocumentCommand{\fnorm}{m}{\ensuremath{\lVert#1\rVert_{\mathrm{F}}}} \NewDocumentCommand{\matA}{}{\mathbf{A}}
\NewDocumentCommand{\matB}{}{\mathbf{B}} \NewDocumentCommand{\matC}{}{\mathbf{C}}
\NewDocumentCommand{\matD}{}{\mathbf{D}} \NewDocumentCommand{\matH}{}{\mathbf{H}}
\NewDocumentCommand{\matI}{}{\mathbf{I}} \NewDocumentCommand{\matP}{}{\mathbf{P}}
\NewDocumentCommand{\matU}{}{\mathbf{U}} \NewDocumentCommand{\matV}{}{\mathbf{V}}
\NewDocumentCommand{\matW}{}{\mathbf{W}} \NewDocumentCommand{\matY}{}{\mathbf{Y}}
\NewDocumentCommand{\matSigma}{}{\boldsymbol{\Sigma}}

% Notation shorthands
\RenewDocumentCommand{\H}{}{\ensuremath{\mathcal{H}}} % Hankel matrix
\NewDocumentCommand{\cMat}{}{\ensuremath{\mathcal{C}}} % controllability matrix
\NewDocumentCommand{\oMat}{}{\ensuremath{\mathcal{O}}} % observability matrix

% ------------------------------------ Title, authors, affiliations ------------------------------------
\title{State-space modelling of\texorpdfstring{\\}{ }head-related transfer functions}

\author[1]{Art J. R. Pelling} \author[1]{Ennes Sarradj} \correspondingauthor{a.pelling@tu-berlin.de}{Art J. R. Pelling
  and E. Sarradj}

\affil[1]{Department of Engineering Acoustics, TU Berlin, Einsteinufer 25, 10587 Berlin, Germany}

\addbibresource{local-bib.bib}

% ------------------------------------
\begin{document}
\maketitle

%%%%%%%%%%%%%%%%%%%%%%%%%%%%%%%%%%%%%%%%%%%%%%%%%%%%%%%%%%%%%%%%%%%%%%%%%%%%%%%%
\begin{abstract}
  State-space models can provide an efficient alternative to convolution-based rendering of head-related transfer
  functions (HRTFs), but directly identifying reduced state-space models from densely sampled datasets can be
  computationally infeasible. We present a loss-free projected Eigensystem Realization Algorithm (ERA) formulation for
  multiple-input multiple-output HRTF datasets, in which sampled directions are inputs and ears are outputs. The
  projection basis preserves the nonzero Hankel singular values exactly and therefore adds no approximation beyond
  dynamical order reduction. For finite impulse-response approximations, we prove that both the McMillan degree and the
  loss-free projection dimension are bounded by the product of the response length and number of outputs, independently
  of the number of sampled directions, and derive a projection-aware ERA error bound. For the 11\,950-input, 256-sample
  FABIAN dataset, the input projection to 510 virtual inputs reduces the memory requirements of the economy-SVD ERA-step
  from 23.2\,GiB to approximately 1.0\,GiB, enabling computation of reduced-order state-space models on virtually any
  platform. We further show that interaural-time-delay removal accelerates singular-value decay and can greatly improve
  quality of identified models.
\end{abstract}

\keywords{model order reduction, balanced truncation, head-related transfer function, auralization, data-driven
  modelling}

%%%%%%%%%%%%%%%%%%%%%%%%%%%%%%%%%%%%%%%%%%%%%%%%%%%%%%%%%%%%%%%%%%%%%%%%%%%%%%%%
%
%
%
\section{INTRODUCTION}
\label{sec:introduction}
%
Immersive virtual acoustic rendering requires listener-specific \acp{HRTF} that remain accurate as source and listener
orientations change \cite{adams2008,moller1992}. Direction-dependent spectral changes and \acp{ITD} support sound source
localization \cite{georgiou1999,moller1992}. \ac{HRTF} models must therefore accommodate real-time source and head
motion with adequate directional resolution \cite{adams2008}.

Virtual acoustic rendering commonly selects directionally sampled \acp{HRIR} in the form of \ac{FIR} filters and
convolves them either in the time or frequency domain, e.g. using the \ac{UPOLS} method \cite{wefers2015}. State-space
models provide a less widespread alternative that can reduce real-time \ac{MIMO} rendering cost, especially after order
reduction \cite{adams2007b,adams2008,adams2009a,pelling2021}.

Balanced truncation and \ac{ERA} have produced reduced order state-space \ac{HRTF} models for individual responses and
multi-direction datasets \cite{mackenzie1997,georgiou1999,grantham2005,adams2007b,adams2008,adams2009a}. They use a truncated
\ac{SVD} of the Hankel matrix, whose cost grows with the number of measured directions and can become problematic for
densely sampled \ac{HRTF} datasets \cite{kramer2016,minster2021}. Tangential projection, CUR factorization, and
randomized \ac{SVD} can reduce this cost by trading approximation quality for efficiency
\cite{kramer2016,kramer2018,minster2021,pelling2021,pelling2026}. Here, we exploit the special structure of \ac{HRTF} datasets to
compute a loss-free projected factorization that exactly preserves the Hankel singular values at a greatly reduced cost.

This paper identifies reduced-order \ac{MIMO} state-space models from densely sampled \acp{HRTF}, with measurement
directions as inputs and ears as outputs. We call \ac{ERA} applied to the full, unprojected Hankel matrix direct
\ac{ERA}; loss-free projected \ac{ERA} instead uses an exact virtual input basis without tangential truncation, reducing
only the model's dynamical behaviour, not its spatial fidelity. We prove bounds on the McMillan degree, loss-free
virtual input dimension, and reduced-model error, and evaluate the method using the FABIAN database
\cite{brinkmann2017}. We also compare original-delay and \ac{ITD}-removed \acp{HRIR} to assess how delay handling
affects singular-value decay and modelling error.
%
%%%%%%%%%%%%%%%%%%%%%%%%%%%%%%%%%%%%%%%%%%%%%%%%%%%%%%%%%%%%%%%%%%%%%%%%%%%%%%%%
%
%
%
\section{PROJECTED \acs{ERA} for \acp{HRTF}}
\label{sec:methods}
%
In the following, we will consider MIMO discrete-time state-space systems with \(m\in\N\) inputs and \(p\in\N\) outputs.
Such a system can be described by a so-called \emph{realization} given by the matrix quadruple
\((\matA,\matB,\matC,\matD)\) with \(\matA\in\R[n][n]\), \(\matB\in\R[n][m]\), \(\matC\in\R[p][n]\), and
\(\matD\in\R[p][m]\), where \(n\in\N\) is the order of the system. Its dynamics are governed by the discrete-time state
equations
%
\begin{align*}
  x_{k+1}&=\matA x_k+\matB u_k,\\y_k&=\matC x_k+\matD u_k,
\end{align*}
%
for a causal input \(u\in\ell_{2}^{m}\), output \(y\in\ell_{2}^{p}\) and state \(x\in\ell_{2}^{n}\), where \(\ell_2^d\)
denotes the square-summable sequences of \(d\)-dimensional vectors and \(k\in\N_0\) is the discrete-time index. Assuming
a zero initial state, a stable realization, and \(u_k=0\) for \(k<0\), the input-output behaviour is fully described by
\(y_k=\sum_{\ell\in\N_0}h_\ell u_{k-\ell}\), where the Markov parameters are
%
\begin{equation}
  h_k=
  \begin{cases}
    \matD, & k=0,\\
    \matC\matA^{k-1}\matB, & k\geq1.
  \end{cases}
  \label{eq:markov-parameters}
\end{equation}
%
Equivalently, in the frequency domain, the output \(\matY(z)\in\C[p]\) is given by a multiplication of the input
\(\matU(z)\in\C[m]\) with the transfer function
%
\begin{equation*}
  \matH(z)=\matC(z\matI_n-\matA)^{-1}\matB+\matD,
\end{equation*}
%
where \(z\) is the frequency and \(\matI_n\in\R^{n\times n}\) is the identity.

For the \ac{HRTF} dataset considered below, each measured source direction is an input and the two ears are outputs.
Thus, one realization shares a state matrix across every direction-ear transfer function, making it a \ac{CAP} model in
the representational sense of \cite{haneda1999a}.
%
%
\subsection{ERA}
\label{sec:era}
%
In discrete time, the impulse response is the sequence of Markov parameters \cref{eq:markov-parameters}. This
relationship allows the identification of a matching state-space model from an impulse response sequence
\((h_0,h_1,\dots,h_{2s-1})\), \(s\in\N\), via the Hankel matrix of Markov parameters
%
\begin{equation}
  \label{eq:H}
  \H=
  \begin{bmatrix}
    \matC\matB & \matC\matA\matB & \cdots & \matC\matA^{s-1}\matB\\
    \matC\matA\matB & \matC\matA^2\matB &\cdots & \matC\matA^s\matB \\
    \vdots & \vdots & \iddots &\vdots\\
    \matC\matA^{s-1}\matB & \matC\matA^{s}\matB & \cdots & \matC\matA^{2s-2}\matB
  \end{bmatrix}.
\end{equation}
%
The Hankel matrix can be factored into the observability and controllability matrices \(\oMat\in\R[ps][n]\) and
\(\cMat\in\R[n][ms]\):
%
\begin{align*}
  \H=
  \underbrace{\begin{bmatrix}
    \matC\\
    \vdots\\
    \matC\matA^{s-1}
  \end{bmatrix}}_{\eqqcolon \oMat}
  \underbrace{
  \begin{bmatrix}
    \matB&\cdots&\matA^{s-1}\matB
  \end{bmatrix}}_{\eqqcolon \cMat}\in\R[ps][ms].
\end{align*}
%
From this factorization, a realization can be constructed via
%
\begin{equation}
  \label{eq:realization}
  \matA=\oMat_{f}^{\dagger}\oMat_{l},~\matB=\cMat
  \begin{bmatrix}
    \matI_{m}\\0
  \end{bmatrix},~\matC=[\matI_p~~0]\oMat,~\matD=h_0,
\end{equation}
%
where \(\oMat_{f}\) and \(\oMat_{l}\) denote the first and last \(p(s-1)\) rows of the observability matrix
\(\oMat\in\R[ps][n]\), respectively, and \(\dagger\) denotes the Moore--Penrose pseudoinverse \cite{kung1978,juang1985}.
Classically, the factorization is computed from a (truncated) \acs{SVD}, i.e., for a truncation order
\(r\leq\min\{ms,\,ps\}\) write the SVD as
%
\begin{equation*}
  \H=
  \begin{bmatrix}
    \matU_{r}&\ast
  \end{bmatrix}
  \begin{bmatrix}
    \matSigma_{r}&0\\0&\ast
  \end{bmatrix}
  \begin{bmatrix}
    \matV_{r}^\top\\\ast
  \end{bmatrix},
\end{equation*}
%
where \(\matU_{r}\in\R[ps][r]\) comprises the first \(r\) left singular vectors as columns,
\(\matSigma_{r}=\diag{\sigma_{1},\,\dots,\,\sigma_{r}}\) the first \(r\) singular values and \(\matV_{r}\in\R[ms][r]\)
the first \(r\) right singular vectors as columns. By choosing
%
\begin{equation}
  \label{eq:grams}
  \oMat=\matU_{r}\matSigma_{r}^{1/2}\quad\text{and}\quad\cMat=\matSigma_{r}^{1/2}\matV_{r}^\top,
\end{equation}
%
the resulting realization is approximately internally balanced \cite{kung1978,schutter2000}. Assuming that the impulse
response has decayed to negligible magnitude by index \(s\), i.e., \(h_k\approx0\) for \(k>s\), the resulting
realization is stable \cite{pelling2026,kung1978,kramer2016}. Additionally, the realization given by
\cref{eq:realization,eq:grams} satisfies the a priori error bound
%
\begin{equation}
  \label{eq:original-bound}
  \sum_{k=1}^{2s-1}\fnorm{\matC_{r}\matA_{r}^{k-1}\matB_{r}-h_{k}}^2\leq(r+m+p)\cdot\sigma_{r+1}(\H)^2,
\end{equation}
%
where \(\sigma_{r+1}(\H)\) is the first neglected Hankel singular value and \(\fnorm{\cdot}\) is the Frobenius norm
\cite{pelling2026,kung1978}.
%
%
\subsection{Loss-Free Projected \acs{ERA}}
\label{sec:tangential}
%
By \cref{sec:era}, we can construct reduced order models from \ac{HRIR} datasets\footnote{\acp{HRTF} with equidistant
  frequency bins can be transformed into \acp{HRIR} with an FFT.}. Recall that \(\H\in\R^{ps\times ms}\). For datasets
with many input directions, \(m\) is large, which can make the block-Hankel matrix infeasibly wide for a direct \ac{SVD}
factorization. A natural approach would be to apply tangential projections with a reduced number of inputs in the
\ac{TERA} framework \cite{kramer2016}. As mentioned, this can accelerate the computations albeit yielding a suboptimal
reduced order model. As we will show now, the special structure of \acp{HRIR} allows us to omit tangential truncation
and use an exact basis for the row space of \(\H\) for projection instead.

\acp{HRIR} typically decay within a few milliseconds, so it is common practice to consider \ac{FIR} approximations where
the \ac{HRIR} is set to zero after some time index \(s\). The following theorem bounds the McMillan degree of an
\ac{FIR} approximation with \(s\) nonzero samples:
%
\begin{theorem}[McMillan degree bound]
  \label{thm:order-bound}
  Let \(h_k\in\R[p][m]\), \(k=0,\dots,s\), denote an \ac{FIR}-approximated impulse response sequence with \(h_k=0\) for
  every \(k>s\). Then, the rational transfer function given by \(\matH_s(z)=\sum_{k=0}^sh_kz^{-k}\) has McMillan degree
  at most \(ps\). In particular, for a \ac{MIMO}-\ac{HRTF}, \(p=2\) (the two ears), so its McMillan degree is at most
  \(2s\), independently of the number \(m\) of sampled directions.
\end{theorem}
%
\begin{proof}
  Let \(\H_\infty=[h_{i+j-1}]_{i,j\geq1}\) be the infinite Hankel matrix of \(\matH_s\). Because the sequence is zero
  for \(k>s\), \(\H_\infty\) is zero outside its leading \(s\times s\) block, which is \(\H\) in \cref{eq:H}. Thus
  \(\rank{\H_\infty}=\rank{\H}\leq ps\). The McMillan degree of a rational function equals the rank of its infinite
  Hankel matrix, proving the result.
\end{proof}
%
This bound is a property of finite \ac{HRIR}, independent of its realization. Thus, it applies whether or not a
state-space model is constructed. Since \cref{eq:H} captures the Hankel singular values of an \(s\)-sample
\ac{MIMO}-\ac{HRIR} entirely, retaining its \(r=\rank{\H}\leq2s\) nonzero singular components in \ac{ERA} exactly
realizes the corresponding zero-extended \ac{HRTF}. Thus, \(2s\) is a worst-case order bound rather than necessarily the
minimal order. Following input-side \ac{TERA} \cite{kramer2016}, let
%
\begin{equation*}
  \Theta_R=\begin{bmatrix}h_1^\top&\cdots&h_s^\top\end{bmatrix}^\top
  =\bar{\matU}\bar{\matSigma}\bar{\matV}^\top\in\R[ps][m]
\end{equation*}
%
be its economy \ac{SVD}\footnote{An economy-size/rank-revealing LQ factorization can also be used here.}. The direct
term is retained without projection, so \(\matD_r=\matD=h_0\). Unlike standard \ac{TERA}, we omit tangential truncation:
with \(q=\rank{\Theta_R}\), set \(\matW_2=\bar{\matV}\in\R[m][q]\), \(\matP_2=\matW_2\matW_2^\top\), and define the
projected Markov parameters \(\hat h_k=h_k\matW_2\).
%
\begin{theorem}[Loss-free virtual inputs]
  \label{thm:virtual inputs}
  Assume \(h_k=0\) for every \(k>s\). Then, \(q\leq\rank{\H}\leq ps\) and \(h_k=h_k\matP_2\) for every \(k\geq1\).
  Consequently, the projected Hankel matrix constructed from \(\hat h_k=h_k\matW_2\) has the same nonzero Hankel
  singular values as \cref{eq:H}. For every reduced order, the Markov parameters of the projected and unprojected
  \acs{ERA} models are identical.
\end{theorem}
%
\begin{proof}
  \(\Theta_R\) is the first block column of \(\H\), and \(\matW_2\) spans \(\operatorname{range}(\Theta_R^\top)\); hence
  \(q\leq\rank{\H}\) and \(h_k=h_k\matP_2\). With \(\mathcal{W}=\matI_s\otimes\matW_2\),
  \(\H=\H\mathcal{W}\mathcal{W}^\top\), so \(\hat{\H}=\H\mathcal{W}\) has the same nonzero singular values.
  Corresponding truncated \acp{SVD} differ only in their input coordinates; thus \cref{eq:grams} gives the same reduced
  Markov parameters after backprojection.
\end{proof}

\Cref{thm:order-bound,thm:virtual inputs} give a projected \ac{ERA} workflow that replaces a large \ac{SVD} of size
\((2s\times ms)\) by two smaller \acp{SVD} of sizes \((2s\times m)\) and \((2s\times qs)\) that can lead to substantial
gains when \(q\ll m\): First, compute \(\matP_2\) from the measured data, identify a reduced realization
\((\matA_r,\hat\matB_r,\matC_r,h_0)\) from \(\hat h\) with \ac{ERA}, then backproject by
\(\matB_r=\hat\matB_r\matW_2^\top\). The factor \(\matW_2\) of \(\matP_2\) maps a physical input direction \(e_j\) to
virtual input coefficients \(\matW_2^\top e_j\), encoding each virtual input as a linear combination of input
directions. For a \acs{HRTF} dataset, \(q\leq2s\), even when additional measured directions are added, so this strategy
can be applied to \ac{HRTF} datasets of practically any level of spatial fidelity. Finally, the previous rank
considerations can be used to sharpen the a priori error bound for \ac{ERA}:
%
\begin{theorem}[Projection-aware error bound]
  \label{thm:projected-bound}
  Under the assumptions of \cref{thm:virtual inputs}, let \((\matA_r,\hat\matB_r,\matC_r,h_0)\) be the order-\(r\)
  \acs{ERA} realization of \(\hat h\). Its backprojected realization \((\matA_r,\hat\matB_r\matW_2^\top,\matC_r,h_0)\)
  satisfies
  \begin{equation*}
    \sum_{k=1}^{2s-1}\fnorm{\matC_r\matA_r^{k-1}\hat{\matB}_r\matW_2^\top-h_k}^2\leq(r+q+p)\cdot\sigma_{r+1}(\H)^2.
  \end{equation*}
\end{theorem}
%
\begin{proof}
  Apply \cref{eq:original-bound} to \(\hat h\), which has \(q\) inputs. By \cref{thm:virtual inputs}, the singular
  values agree and backprojection preserves the error norm because \(\matW_2\) has orthonormal columns.
\end{proof}

%%%%%%%%%%%%%%%%%%%%%%%%%%%%%%%%%%%%%%%%%%%%%%%%%%%%%%%%%%%%%%%%%%%%%%%%%%%%%%%%
%
%
%
\section{EXPERIMENTAL SETUP}
\label{sec:experimental-setup}
%
%
\subsection{FABIAN Dataset and Preprocessing}
\label{sec:fabian}
We use the \(256\)-sample, \(44\,100\,\)Hz \acp{HRIR} for the neutral \acs{HATO} in the FABIAN \acs{HRTF} database
\cite{brinkmann2017} with \(m=11\,950\) source directions at \(2^{\circ}\) elevation and azimuth increments. We treat
them as \(h_0,\ldots,h_{255}\) of a \ac{MIMO} discrete-time system with \(p=2\) outputs (the two ears), so \(s=255\)
non-direct Markov parameters are \(h_1,\ldots,h_{255}\). We zero-extend the sequence by \(h_k=0\) for \(k>s\). Let
\(h^{(i,j)}\in\R^{2s}\) denote the resulting HRIR from the \(i\)-th source to the \(j\)-th receiver. The \(k\)-th Markov
parameter contains the \(k\)-th sample of every HRIR, i.e. \(\{h_k\}_{i,j}=h^{(i,j)}_k\in\R[p][m]\). Before system
identification, the \acp{HRIR} are normalized so that the mean energy of single-channel \acp{HRTF} is one.

%
%
\subsection{Delay Handling}
\label{sec:delay-handling}
We use two delay-handling conditions. \emph{Original-delay} \acp{HRIR} are normalized, zero-extended responses
without shifts. For \emph{\ac{ITD}-removed} \acp{HRIR}, we estimate each response onset as the last sample below
\SI{20}{\decibel} relative to its maximum and shift the response by this estimate\footnote{See
  \href{https://pyfar.readthedocs.io/en/v0.8.0/modules/pyfar.dsp.html\#pyfar.dsp.find_impulse_response_start}{\texttt{find\_impulse\_response\_start}
    (pyFAR v0.8.0)}.}. For the FABIAN dataset, this retains \(s=\ITDRemovedLength\) nonzero samples.
%
\subsection{Implementation and Timing}
\label{sec:timing}
Timings cover \ac{HRIR} preprocessing, projected \ac{ERA} factorization, and reduced-realization construction, excluding
data loading. Three uncached trials per case ran on an AMD Ryzen~7 8845HS laptop (8 cores, 16 threads; \(14\)\,GiB
memory) under Arch Linux, Python~3.13, NumPy~2.4.6, and SciPy~1.18.0 with 16 OpenBLAS threads.

%%%%%%%%%%%%%%%%%%%%%%%%%%%%%%%%%%%%%%%%%%%%%%%%%%%%%%%%%%%%%%%%%%%%%%%%%%%%%%%%
%
%
%
\section{RESULTS}
\label{sec:results}
%
\begin{figure*}[tb]
  \centering \setlength{\hsvmapheight}{0.263\textwidth}
  \begin{subfigure}[t]{0.42\textwidth}
    \centering \includegraphics[height=\hsvmapheight]{../generated/figures/hsv-map.pdf}
    \subcaption{Original-delay \acp{HRIR}; colors show singular values.}
    \label{fig:hsv-map}
  \end{subfigure}\hspace{0\textwidth}
  \begin{subfigure}[t]{0.42\textwidth}
    \centering \includegraphics[height=\hsvmapheight]{../generated/figures/hsv-map-itd-removed.pdf}
    \subcaption{After individual \ac{ITD} removal.}
    \label{fig:hsv-map-itd-removed}
  \end{subfigure}\hspace{-0.02\textwidth}
  \begin{subfigure}[t]{0.10\textwidth}
    \centering \includegraphics[height=\hsvmapheight]{../generated/figures/hsv-colorbar.pdf}
  \end{subfigure}
  \caption{Hankel singular-value maps for neutral-orientation FABIAN \acp{HRIR}. Colors are dB relative to the maximum;
    red lines mark loss-free virtual input caps and circles \protect\(r=\SelectedModelOrder\).}
  \label{fig:hsv-maps}
\end{figure*}
%
\subsection{Loss-Free Virtual Input Projection}
\label{sec:projection-results}
For the original-delay FABIAN responses, the number of loss-free virtual inputs is \(\OriginalDelayRankBound\); after
individual \ac{ITD} removal it is \(2\cdot\ITDRemovedLength=\ITDRemovedRankBound\). \Cref{fig:hsv-maps} depicts the Hankel singular values
decay for a varying (truncating) number of virtual inputs in the \ac{TERA} framework. The red vertical lines mark the
respective values for \(q\) based on \cref{thm:virtual inputs} at which the input projection becomes loss-free. To the
right of these, the contours remain horizontal, showing that further virtual inputs leave the Hankel singular values
unchanged. Several contours are already nearly horizontal to the left of these lines, suggesting that a lossy tangential
truncation could further reduce the input dimension for a chosen accuracy target. We stick to \(q\) based on
\cref{thm:virtual inputs} because the computational gains are already sufficient and an exact preservation of the
\acp{HRTF} is preferable.

\Cref{fig:scalability} shows that direct economy-SVD storage would require \(\DirectSVDMemoryGiB\)\,GiB. Capping
the virtual input dimension at \(q=\OriginalDelayRankBound\) reduces storage to approximately
\(\ProjectedSVDMemoryGiB\)\,GiB and full-SVD work by a factor of about \(\SVDWorkReductionFactor\). These savings
are theoretical, not measured speedups, because the direct factorization exceeded the available memory of our platform.
%
\begin{figure}[tb]
  \centering
  \includegraphics[width=\linewidth]{../generated/figures/scalability.pdf}\\[-0.15\baselineskip]
  \includegraphics[width=\linewidth]{../generated/figures/scalability-legend.pdf}\\[-0.2\baselineskip]
  \setlength{\abovecaptionskip}{1pt}
  \caption{Economy-SVD memory and full-SVD work versus sampled directions, including dense Hankel and economy-size
    factors. The red line marks the loss-free virtual input cap \protect\(q=\OriginalDelayRankBound\).}
  \label{fig:scalability}
\end{figure}
%
\begin{figure}[tb]
  \centering
  \includegraphics[width=\linewidth]{../generated/figures/bound-error.pdf}\\[-0.15\baselineskip]
  \includegraphics[width=\linewidth]{../generated/figures/bound-error-legend.pdf}\\[-0.2\baselineskip]
  \setlength{\abovecaptionskip}{1pt}
  \caption{Hankel singular-value decay, a priori bounds, and relative \acs{HRIR} reconstruction errors for
    neutral-orientation FABIAN data. Colors distinguish delay handling; line styles distinguish quantities; and circles
    mark restored-response errors at \protect\(r=\SelectedModelOrder\); the dark curve uses the aligned reference. For
    \ac{ITD}-removed responses, the bounds apply to the aligned system, so comparison with restored-response errors is
    empirical.}
  \label{fig:bound}
\end{figure}
%
\subsection{Effect of \acs{ITD}-Alignment}
\label{sec:itd-results}
%
\Cref{fig:bound} compares singular-value decay, the two a priori bounds, and measured relative \ac{HRIR} reconstruction
errors. As can be seen from \cref{fig:hsv-maps,fig:bound}, individual \ac{ITD} removal produces a faster singular value
decay and reduces the relative error at a fixed reduced order. However, the relative error of \ac{ITD}-restored reduced
order models does not decrease below \SI{\ITDRestoredHighOrderErrorDB}{\decibel}. At \(r=\ITDRestoredCrossoverOrder\),
the \ac{ITD}-restored models perform worse than the original-delay models. A further increase of model order does not
improve the error substantially. This, however, does not stem from the projected \ac{ERA} routine but is due to the
truncation of nonzero values of the \ac{HRIR} when removing the estimated \acp{ITD}. The \ac{ITD}-aligned error also
depicted in \cref{fig:bound} continues to decay to \SI{\ITDAlignedHighOrderErrorDB}{\decibel} at order
\(\HighestEvaluatedOrder\), showing that the high-order plateau is not caused by a lack of convergence
against the aligned reference. An increase in model order does in fact improve the approximation of the \ac{ITD}-removed
\ac{HRIR} but at a certain point, the error in truncation during removal becomes larger than the error in the model
reduction.

Earlier \ac{HRTF} work likewise found that alignment or phase correction can reduce effective spatial order
\cite{ben-hur2019a}, and that \ac{ITD} can degrade reduced state-space models \cite{adams2008,pelling2026}. The present
result is consistent with these observations, but it does not establish a causal mechanism.
%
\subsection{Order Selection}
\label{sec:response-results}
%
The common reduced order \(r=\SelectedModelOrder\) was selected as the lowest order whose original-delay reconstruction
error is below \SI{-12}{\decibel} (\SI{\OriginalDelayRelativeErrorDB}{\decibel}). This threshold is set arbitrarily and not directly motivated by a perceptual study.
Earlier \ac{HRTF} state-space studies reported and perceptual order thresholds
between 7 and 24 for an array of 50 \acp{HRTF} \cite{adams2009a}, but these conditions are not directly comparable with
the present 11\,950-direction dataset. The literature therefore only contextualizes the order scale; perceptual and
localization validation is still required.

\Cref{fig:responses} examines the individual left and right responses at azimuth \SI{45}{\degree} and elevation
\SI{15}{\degree}. The time-domain panel compares the measured responses with the delay-restored, \ac{ITD}-removed
order-\(\SelectedModelOrder\) model. The magnitude panel compares the measured responses, the original-delay and
\ac{ITD}-removed models, and the corresponding ear-specific errors. Thus, the individual-response comparison
complements the aggregate errors in \cref{fig:bound} at the selected common order.
%
\begin{figure*}[tb]
  \centering
  \begin{subfigure}[t]{0.49\textwidth}
    \centering \includegraphics[width=\linewidth]{../generated/figures/responses-itd-removed-time.pdf}
    \subcaption{Individual \ac{ITD} removal, time domain.}
  \end{subfigure}\hfill
  \begin{subfigure}[t]{0.49\textwidth}
    \centering \includegraphics[width=\linewidth]{../generated/figures/responses-frequency.pdf}
    \subcaption{Both delay-handling methods, magnitude domain.}
  \end{subfigure}\\[-0.25\baselineskip]
  \begin{subfigure}[t]{\textwidth}
    \centering \includegraphics[width=\linewidth]{../generated/figures/responses-legend.pdf}
  \end{subfigure}\\[-0.3\baselineskip]
  \setlength{\abovecaptionskip}{0pt}
  \caption{Measured and order-\protect\(\SelectedModelOrder\) FABIAN responses at azimuth \SI{45}{\degree} and elevation
    \SI{15}{\degree}: (a) the individually \ac{ITD}-removed model in time and (b) both delay-handling methods in
    magnitude. Black curves are measurements; normal and lighter shades identify reduced models and errors; dashed
    blue/orange curves are original-delay left/right results.}
  \label{fig:responses}
\end{figure*}
%
%%%%%%%%%%%%%%%%%%%%%%%%%%%%%%%%%%%%%%%%%%%%%%%%%%%%%%%%%%%%%%%%%%%%%%%%%%%%%%%%
%
%
%
\section{CONCLUSIONS}
\label{sec:conclusions}
%
We presented a projected \ac{ERA} formulation for densely sampled \ac{MIMO} \acp{HRTF}. For finite length-\(s\) \ac{FIR}-\acp{HRIR}, the McMillan degree and loss-free virtual input dimension are bounded by \(2s\), independently of the number
of sampled directions. This data-driven bound preserves the nonzero Hankel singular values and can also guide
investigation of other Hankel- or \ac{SVD}-based \ac{HRTF} representations.

For the \(11\,950\)-direction FABIAN dataset, loss-free projection reduces estimated economy SVD storage from \(\DirectSVDMemoryGiB\)\,GiB to approximately
\(\ProjectedSVDMemoryGiB\)\,GiB, lifting the memory limit of the evaluated laptop platform. The projection-aware error
bound is correspondingly tighter than the classical \acs{ERA} bound because it uses \(q=510\) rather than all \(m=11\,950\) input
directions. Individual \ac{ITD} removal accelerated singular-value decay and reduced the relative
reconstruction error from \SI{\OriginalDelayRelativeErrorDB}{\decibel} to \SI{\ITDRemovedRelativeErrorDB}{\decibel}.
The study considers one neutral \ac{HATO} of the FABIAN head. Future work should test further \acp{HATO} and directional
samplings, assess perceptual and localization implications, and evaluate dead-time-splitting extraction
\cite{pelling2026} for improved real-time rendering.

%%%%%%%%%%%%%%%%%%%%%%%%%%%%%%%%%%%%%%%%%%%%%%%%%%%%%%%%%%%%%%%%%%%%%%%%%%%%%%%%
\section*{Source Code Availability}
The source code used to compute the results is available at
\href{https://github.com/artpelling/fa26-hrtf-ssm}{\texttt{GitHub:artpelling/fa26-hrtf-ssm}} under the MIT license.
%
\section*{Acknowledgements}
The work of Art J. R. Pelling was funded by DFG (German Research Foundation) – project no.:
\href{https://gepris.dfg.de/gepris/projekt/504367810?language=en}{504367810}.
%
\begingroup
% \renewcommand*{\bibfont}{\small}
\printbibliography[title={References}] \endgroup
%
\end{document}
